\section{Oil for the Lamps, Fruit from the Talents}

The two parables preceding the final judgment in Matthew~25 teach readiness through the passing of time. The bridegroom delays; the master returns after a long absence. The virgins must have oil when the summons comes, and the servants must account for their use of what the master entrusted to them. Preparation takes place during the absence. The eventual return makes their conduct during that time decisive.

All ten virgins fall asleep. The difference lies in their provision of oil. When the summons comes, the wise are ready to enter; the foolish have left the necessary preparation undone (Mt 25:1--13). Vigilance therefore cannot mean a state of uninterrupted emotional intensity. Human beings must sleep, work, eat, and fulfill ordinary responsibilities. The readiness Christ commands can inhabit these things. It consists in a life furnished for his coming, one in which the inward substance of love accompanies the outward profession of expectation.

Chrysostom understands the lamps as the honor of virginity and the oil as mercy, especially almsgiving. The great labor of chastity does not dispense the virgins from love of the needy (LXXVIII.1). This interpretation bears on other forms of religious accomplishment. We may persevere in a difficult penance while refusing an ordinary duty of care. The good of the penance remains a reason for thanksgiving, but the refusal still requires repentance. Examination must reach the duty we have evaded even when another part of our life has demanded considerable sacrifice.

The foolish virgins discover their lack when the opportunity to remedy it has passed. In Chrysostom's exposition, the request for borrowed oil shows that another person's goodness cannot replace the response which should have been ours (LXXVIII.1--2). The communion of saints truly helps us, and their intercession is a gift; it does not turn our refusal to love into their obedience. We cannot borrow a parent's fidelity while rejecting the faith the parent taught, or rest upon the holiness of a community while refusing its discipline. Grace reaches us through others so that we too may respond.

The parable of the talents deepens the same summons. The master entrusts different amounts according to the servants' capacities. The servants who receive five and two both put the gift to work, and both receive the same commendation for fidelity. The servant who receives one buries it (Mt 25:14--30). Christ asks for a faithful answer to what has actually been entrusted. A smaller capacity does not prevent that answer. Envy of another person's gifts can become an excuse for leaving our own unused.

Chrysostom names a wide range of talents: money, teaching, assistance, protection, and whatever ability a person can use for the neighbor's good (LXXVIII.2--3). A familiar ability may be useful to someone who lacks it. Reading a difficult letter for a neighbor can help him understand a responsibility he otherwise could not fulfill. The task requires no unusual public achievement, yet refusing it out of indifference would leave a capacity unused when it was needed. Chrysostom's examples invite us to examine what we can offer.

The buried talent is especially sobering because the servant returns what he received. He can point to its preservation, yet he has refused the purpose of the gift. His account of the master as harsh becomes part of the excuse. Fear has protected inactivity. This can happen in spiritual life when we use our imperfection as a reason never to begin serving, or the possibility of failure as a reason never to attempt a known good. We may be unable to guarantee success; we remain able to obey within our responsibility.

Chrysostom applies this distinction to instruction and correction. The servant can give the counsel that is his to give and leave the outcome to the Lord; another person's refusal does not make the faithful effort worthless (LXXVIII.2). A correction can be necessary even when it is unwelcome. After giving it with truth and charity, the person responsible may have to wait, remain available, and continue the other duties of care. The hearer's refusal remains the hearer's choice.

The warning against delay continues in the preceding servant parable. A servant decides that his master is long in returning, begins to beat his fellow servants, and indulges himself. The delay becomes permission for abuse (Mt 24:45--51). In contrast, the faithful servant gives the household its food at the appointed time. Readiness is shown in fulfilling the duty for which authority was given. Parents, employers, teachers, and those entrusted with any responsibility prepare for judgment by serving those under their care, especially when no superior is watching.

The passing years can change the way we fulfill a responsibility. Familiarity may make patient service easier, or allow resentment to become habitual. Success may make us grateful, or less willing to accept correction. Annual examination can compare our present conduct with the duty itself. A practice once begun faithfully may need renewal; a failure once confessed may have returned under a new excuse. Readiness requires continuing fidelity while the master is absent.

